\documentclass[10pt]{extarticle}
\usepackage[a4paper,margin=1.5cm]{geometry}
\usepackage[T1]{fontenc}
\usepackage{amsmath,amssymb}
\usepackage{graphicx}
\usepackage[export]{adjustbox}
\usepackage{multicol}
\usepackage{enumitem}
\usepackage{booktabs,array,tabularx}
\setlength{\columnsep}{0.7cm}
\setlength{\parindent}{0pt}
\setlength{\parskip}{3pt}
\pagestyle{empty}
\makeatletter
% figure registry (figdefs.tex is generated by the extract step)
\newcommand{\qfigdef}[4]{\@namedef{qf@file@#1}{#2}\@namedef{qf@w@#1}{#3}\@namedef{qf@wide@#1}{#4}}
\newcommand{\qf@check}[1]{\@ifundefined{qf@file@#1}{\PackageError{score2}{Unknown figure id #1}{}}{}}
\newcommand{\qf@img}[4][]{\includegraphics[width=#2\dimexpr\@nameuse{qf@w@#4}pt\relax,max width=#3,max height=0.45\textheight,#1]{\@nameuse{qf@file@#4}}}
% Full-width questions: finalize wraps every question that contains a wide figure in \qfullbegin...\qfullend,
% so the WHOLE question (stem, figures, options) is set across the page in reading order, then the two columns
% resume. The question's head - its number, stem and first figure (\qfullsplit marks the end of the head) -
% is kept on one page, so a number is never left at a page foot with its figure overleaf; the rest may break
% across pages normally, so no big gaps are left. Without a \qfullsplit, the whole question is kept together
% if it is under 45% of a page. Never write these yourself.
\newif\ifqf@full
\newif\ifqf@boxed
\newsavebox{\qf@box}
\newcommand{\qf@open}{\global\qf@fulltrue\par\global\qf@boxedtrue
  \setbox\qf@box\vbox\bgroup\hsize\textwidth\linewidth\textwidth}
\newcommand{\qf@emit}[1]{\par\egroup\global\qf@boxedfalse
  \ifdim\dimexpr\ht\qf@box+\dp\qf@box\relax<#1\textheight\noindent\usebox{\qf@box}\par\else\unvbox\qf@box\fi}
\newcommand{\qfullsplit}{\ifqf@boxed\qf@emit{0.85}\fi}
\newcommand{\qf@close}{\ifqf@boxed\qf@emit{0.45}\fi\global\qf@fullfalse}
\newcommand{\qfullbegin}{\end{multicols}\qf@open}
\newcommand{\qfullend}{\qf@close\begin{multicols}{2}}
\newcommand{\qfullstart}{\qf@open}
\newcommand{\qfulldone}{\qf@close}
% \qfig[scale]{id}: block figure, centred, directly where it is written (size set by the tool).
% Fallback if a wide figure is not inside a full-width question: break out of the columns in place - no floats.
\newcommand{\qfig}[2][1]{\qf@check{#2}%
  \ifqf@full\qf@blk{#1}{#2}\else
  \ifnum\@nameuse{qf@wide@#2}=1 \expandafter\@firstoftwo\else\expandafter\@secondoftwo\fi
  {\end{multicols}\begin{center}\qf@img{#1}{\textwidth}{#2}\end{center}\begin{multicols}{2}}%
  {\qf@blk{#1}{#2}}\fi}
\newcommand{\qf@blk}[2]{\par\smallskip{\centering\qf@img{#1}{\linewidth}{#2}\par}\smallskip}
% \qfigi[scale]{id}: inline figure (inside option lists / tables / minipages); never breaks columns.
\newcommand{\qfigi}[2][1]{\qf@check{#2}\qf@img[valign=c]{#1}{\linewidth}{#2}}
\makeatother
% structure macros - lint reads these to check completeness and figure placement
\newcommand{\papersec}[1]{\par\bigskip{\centering\Large\bfseries PAPER-#1\par}\medskip}
\newcommand{\partsec}[2]{\par\medskip\noindent{\large\bfseries PART-#1: #2}\par\smallskip}
\newcommand{\sect}[1]{\par\smallskip\noindent\textbf{#1}\par}
\newcommand{\qn}[1]{\par\smallskip\noindent\textbf{#1.}~}
\newcommand{\paraq}[2]{\par\smallskip\noindent\textbf{Paragraph for Questions #1 and #2}\par}
\newenvironment{qoptions}{\begin{itemize}[leftmargin=1.8em,itemsep=1pt,topsep=2pt,label={}]}{\end{itemize}}
\newcommand{\opt}[2]{\item[(#1)] #2}
% figure registry (generated by score2_tools.py finalize)
\qfigdef{p1_phy_s1_q01_1}{Star Review (Other) - Test 03 (04 Dec 2023)-img/p1_phy_s1_q01_1.png}{113.4}{0}
\qfigdef{p1_phy_s1_q03_1}{Star Review (Other) - Test 03 (04 Dec 2023)-img/p1_phy_s1_q03_1.png}{108.3}{0}
\qfigdef{p1_phy_s1_q04_1}{Star Review (Other) - Test 03 (04 Dec 2023)-img/p1_phy_s1_q04_1.png}{108.3}{0}
\qfigdef{p1_phy_s2_q06_1}{Star Review (Other) - Test 03 (04 Dec 2023)-img/p1_phy_s2_q06_1.png}{127.5}{0}
\qfigdef{p1_phy_s2_q07_1}{Star Review (Other) - Test 03 (04 Dec 2023)-img/p1_phy_s2_q07_1.png}{139.6}{0}
\qfigdef{p1_phy_s2_q08_1}{Star Review (Other) - Test 03 (04 Dec 2023)-img/p1_phy_s2_q08_1.png}{140.1}{0}
\qfigdef{p1_phy_s2_q10_1}{Star Review (Other) - Test 03 (04 Dec 2023)-img/p1_phy_s2_q10_1.png}{145.4}{0}
\qfigdef{p1_phy_s3_para5_6_1}{Star Review (Other) - Test 03 (04 Dec 2023)-img/p1_phy_s3_para5_6_1.png}{82.0}{0}
\qfigdef{p1_chem_s2_q06_1}{Star Review (Other) - Test 03 (04 Dec 2023)-img/p1_chem_s2_q06_1.png}{137.7}{0}
\qfigdef{p1_chem_s2_q07_1}{Star Review (Other) - Test 03 (04 Dec 2023)-img/p1_chem_s2_q07_1.png}{165.3}{0}
\qfigdef{p1_chem_s2_q07_A}{Star Review (Other) - Test 03 (04 Dec 2023)-img/p1_chem_s2_q07_A.png}{49.9}{0}
\qfigdef{p1_chem_s2_q07_B}{Star Review (Other) - Test 03 (04 Dec 2023)-img/p1_chem_s2_q07_B.png}{51.5}{0}
\qfigdef{p1_chem_s2_q07_C}{Star Review (Other) - Test 03 (04 Dec 2023)-img/p1_chem_s2_q07_C.png}{51.6}{0}
\qfigdef{p1_chem_s2_q07_D}{Star Review (Other) - Test 03 (04 Dec 2023)-img/p1_chem_s2_q07_D.png}{55.7}{0}
\qfigdef{p1_chem_s2_q08_A}{Star Review (Other) - Test 03 (04 Dec 2023)-img/p1_chem_s2_q08_A.png}{90.4}{0}
\qfigdef{p1_chem_s2_q08_B}{Star Review (Other) - Test 03 (04 Dec 2023)-img/p1_chem_s2_q08_B.png}{119.5}{0}
\qfigdef{p1_chem_s2_q08_3}{Star Review (Other) - Test 03 (04 Dec 2023)-img/p1_chem_s2_q08_3.png}{24.9}{0}
\qfigdef{p1_chem_s2_q08_D}{Star Review (Other) - Test 03 (04 Dec 2023)-img/p1_chem_s2_q08_D.png}{54.6}{0}
\qfigdef{p1_chem_s3_q03_1}{Star Review (Other) - Test 03 (04 Dec 2023)-img/p1_chem_s3_q03_1.png}{96.0}{0}
\qfigdef{p1_chem_s3_para5_6_1}{Star Review (Other) - Test 03 (04 Dec 2023)-img/p1_chem_s3_para5_6_1.png}{108.0}{0}
\qfigdef{p2_phy_s1_q03_1}{Star Review (Other) - Test 03 (04 Dec 2023)-img/p2_phy_s1_q03_1.png}{116.1}{0}
\qfigdef{p2_phy_s1_q04_1}{Star Review (Other) - Test 03 (04 Dec 2023)-img/p2_phy_s1_q04_1.png}{129.3}{0}
\qfigdef{p2_phy_s2_para7_8_1}{Star Review (Other) - Test 03 (04 Dec 2023)-img/p2_phy_s2_para7_8_1.png}{101.5}{0}
\qfigdef{p2_phy_s2_para9_10_1}{Star Review (Other) - Test 03 (04 Dec 2023)-img/p2_phy_s2_para9_10_1.png}{129.4}{0}
\qfigdef{p2_phy_s3_para1_2_1}{Star Review (Other) - Test 03 (04 Dec 2023)-img/p2_phy_s3_para1_2_1.png}{103.1}{0}
\qfigdef{p2_phy_s3_para3_4_1}{Star Review (Other) - Test 03 (04 Dec 2023)-img/p2_phy_s3_para3_4_1.png}{148.9}{0}
\qfigdef{p2_phy_s3_para5_6_1}{Star Review (Other) - Test 03 (04 Dec 2023)-img/p2_phy_s3_para5_6_1.png}{130.5}{0}
\qfigdef{p2_chem_s1_q01_1}{Star Review (Other) - Test 03 (04 Dec 2023)-img/p2_chem_s1_q01_1.png}{235.6}{0}
\qfigdef{p2_chem_s1_q02_1}{Star Review (Other) - Test 03 (04 Dec 2023)-img/p2_chem_s1_q02_1.png}{83.5}{0}
\qfigdef{p2_chem_s2_para7_8_4}{Star Review (Other) - Test 03 (04 Dec 2023)-img/p2_chem_s2_para7_8_4.png}{178.7}{0}
\qfigdef{p2_chem_s2_q07_A}{Star Review (Other) - Test 03 (04 Dec 2023)-img/p2_chem_s2_q07_A.png}{42.4}{0}
\qfigdef{p2_chem_s2_q07_B}{Star Review (Other) - Test 03 (04 Dec 2023)-img/p2_chem_s2_q07_B.png}{59.3}{0}
\qfigdef{p2_chem_s2_q07_C}{Star Review (Other) - Test 03 (04 Dec 2023)-img/p2_chem_s2_q07_C.png}{90.3}{0}
\qfigdef{p2_chem_s2_q07_D}{Star Review (Other) - Test 03 (04 Dec 2023)-img/p2_chem_s2_q07_D.png}{63.0}{0}
\qfigdef{p2_chem_s2_q08_A}{Star Review (Other) - Test 03 (04 Dec 2023)-img/p2_chem_s2_q08_A.png}{37.7}{0}
\qfigdef{p2_chem_s2_q08_B}{Star Review (Other) - Test 03 (04 Dec 2023)-img/p2_chem_s2_q08_B.png}{39.2}{0}
\qfigdef{p2_chem_s2_q08_C}{Star Review (Other) - Test 03 (04 Dec 2023)-img/p2_chem_s2_q08_C.png}{41.4}{0}
\qfigdef{p2_chem_s2_q08_D}{Star Review (Other) - Test 03 (04 Dec 2023)-img/p2_chem_s2_q08_D.png}{26.2}{0}
\qfigdef{p2_chem_s3_para3_4_1}{Star Review (Other) - Test 03 (04 Dec 2023)-img/p2_chem_s3_para3_4_1.png}{230.7}{0}
\qfigdef{p2_chem_s4_q07_1}{Star Review (Other) - Test 03 (04 Dec 2023)-img/p2_chem_s4_q07_1.png}{56.1}{0}
\raggedcolumns
\makeatletter\@beginparpenalty=10000\makeatother
\newcolumntype{L}[1]{>{\raggedright\arraybackslash}p{#1}}
\begin{document}
\begin{center}{\Large\bfseries Star Review (Other) --- Test 03 (04 Dec 2023)}\end{center}
\begin{multicols}{2}
\raggedcolumns
\papersec{1}
\partsec{1}{PHYSICS}
\sect{SECTION-I(i)}
\qn{1} In wheat stone's bridge shown in the adjoining figure galvanometer gives no deflection on pressing the key, the balance condition for the bridge is :-\par
\qfig{p1_phy_s1_q01_1}
\begin{qoptions}
\opt{A}{$\dfrac{R_1}{R_2}=\dfrac{C_1}{C_2}$}
\opt{B}{$\dfrac{R_1}{R_2}=\dfrac{C_2}{C_1}$}
\opt{C}{$\dfrac{R_1}{R_1+R_2}=\dfrac{C_2}{C_1-C_2}$}
\opt{D}{$\dfrac{R_1}{R_1-R_2}=\dfrac{C_2}{C_1+C_2}$}
\end{qoptions}
\qn{2} The fractional change in the value of free--fall acceleration $g$ for a particle when it is lifted from the surface to an elevation $h$ ($h\ll R$) is
\begin{qoptions}
\opt{A}{$\dfrac hR$}
\opt{B}{$\dfrac{2h}{R}$}
\opt{C}{$-\dfrac{2h}{R}$}
\opt{D}{$-\dfrac hR$}
\end{qoptions}
\qn{3} A metallic rod of length $\ell$ is rotated in a horizontal plane with angular velocity $\omega$ where no magnetic field exists. The EMF induced in the rod is :- (mass of $e^-$ is $m$ and charge of electron is $e$)\par
\qfig{p1_phy_s1_q03_1}
\begin{qoptions}
\opt{A}{$\dfrac{m\ell\omega^2}{e}$}
\opt{B}{$\dfrac{m\omega\ell^2}{2e}$}
\opt{C}{$\dfrac{m\omega^2\ell^2}{2e}$}
\opt{D}{$\dfrac{m\omega^2\ell^2}{e}$}
\end{qoptions}
\qn{4} A semi-infinite charged wire with charge density $+\lambda_1$ is placed along the axis of a uniformly charged ring of charge density $+\lambda_2$. Find the force exerted on the wire due to the ring. The radius of ring is $R$.\par
\qfig{p1_phy_s1_q04_1}
\begin{qoptions}
\opt{A}{$(k\lambda_1\lambda_2)\,2\pi R$}
\opt{B}{$\dfrac{\lambda_1\lambda_2}{2\epsilon_0}$}
\opt{C}{$k\lambda_1^2\,2\pi R$}
\opt{D}{None of these}
\end{qoptions}
\sect{SECTION-I(ii)}
\qn{5} A $20$ gm particle is subjected to two simple harmonic motions $x_1=2\sin10t$, $x_2=4\sin\left(10t+\dfrac\pi3\right)$. Where $x_1$ and $x_2$ are in metre and $t$ is in sec.
\begin{qoptions}
\opt{A}{The displacement of the particle at $t=0$ will be $2\sqrt3$ m}
\opt{B}{Maximum speed of the particle will be $20\sqrt7$ m/s}
\opt{C}{Magnitude of maximum acceleration of the particle will be $200\sqrt7\ \mathrm{m/s^2}$.}
\opt{D}{Energy of the resultant motion will be $28$ J.}
\end{qoptions}
\qn{6} For the AC circuit shown, the reading of ammeter and voltmeter are $5$A and $50\sqrt5$ volts respectively, then :\par
\qfig{p1_phy_s2_q06_1}
\begin{qoptions}
\opt{A}{average power delivered by the source is $250$W}
\opt{B}{rms value of AC source is $50$ volts}
\opt{C}{rms value of AC source is $100$ volts}
\opt{D}{frequency of ac source is $1000/2\pi$ Hz}
\end{qoptions}
\qn{7} A light ray is incident on a boundary as shown.\par
\qfig{p1_phy_s2_q07_1}
\begin{qoptions}
\opt{A}{If $\mu_2<\sqrt3$, light will enter medium-3 with deviation $30^\circ$.}
\opt{B}{If $\mu_2>\sqrt3$, light will enter medium-3 with deviation $30^\circ$.}
\opt{C}{If $\mu_2>\sqrt3$, light will never return to medium-1 again.}
\opt{D}{If $\mu_2<\sqrt3$, light will never return to medium-1.}
\end{qoptions}
\qn{8} A spring of spring constant $k$ is fixed at one end has a small block of mass $m$ and charge $q$ is attached at the other end. The block rests over a smooth horizontal surface. A uniform and constant magnetic field $B$ exists normal to the plane of paper as shown in figure. An electric field $\vec E=E_0\hat i$ is switched on at $t=0$. The block moves on the horizontal surface without lifting off the surface. Then the normal reaction acting on the the block is :-\par
\qfig{p1_phy_s2_q08_1}
\begin{qoptions}
\opt{A}{Maximum at the mean position.}
\opt{B}{Minimum at the mean position.}
\opt{C}{Equal to $mg$ at extreme position.}
\opt{D}{More than $mg$ at extreme position.}
\end{qoptions}
\qn{9} A neutron collides head-on with a stationary hydrogen atom in ground state. Which of the following statements are correct ? (Assume that the hydrogen atom and neutron has same mass) :-
\begin{qoptions}
\opt{A}{If kinetic energy of the neutron is less than $20.4$ eV collision must be elastic.}
\opt{B}{If kinetic energy of the neutron is less than $20.4$ eV collision may be inelastic.}
\opt{C}{Inelastic collision may be take place only when initial kinetic energy of neutron is greater than $20.4$ eV.}
\opt{D}{Perfectly inelastic collision can not take place.}
\end{qoptions}
\qn{10} In the figure below, the switches $S_1$ and $S_2$ are closed simultaneously at $t=0$ and a current starts to flow in the circuit. Both the batteries have the same magnitude of the electromotive force (emf) and the polarities are as indicated in the figure. Ignore mutual inductance between the inductors. The current $I$ in the middle wire reaches its maximum magnitude $I_{\max}$ at time $t=\tau$. Which of the following statement(s) is (are) true?\par
\qfig{p1_phy_s2_q10_1}
\begin{qoptions}
\opt{A}{$I_{\max}=\dfrac{V}{2R}$}
\opt{B}{$I_{\max}=\dfrac{V}{4R}$}
\opt{C}{$\tau=\dfrac LR\ln2$}
\opt{D}{$\tau=\dfrac{2L}{R}\ln2$}
\end{qoptions}
\sect{SECTION-II(i)}
\paraq{1}{2}
\textbf{Stem-I}\par
A sample of $\mathrm{He^+}$ ions is excited by photons of energy $48.36\ eV$. Photons are emitted when the ions return to their ground state. This radiation falls on a metallic plate of work function $6eV$. [$h=6.6\times10^{-34}$J.s, $m_e=9.1\times10^{-31}$kg]
\qn{1} The maximum wavelength of photon emitted when the $\mathrm{He^+}$ ions return to their ground state is approximately \underline{\hspace{1.5cm}} nm.
\qn{2} The minimum de-Broglie wavelength of the photoelectron emitted is approximately \underline{\hspace{1.5cm}} \AA.
\paraq{3}{4}
${}^{64}_{29}$Cu can decay by $\beta^-$ or $\beta^+$ emission, or electron capture. It is known that ${}^{64}_{29}$Cu has a half life of $12.8$ hrs with $40\%$ probability of $\beta^-$ decay, $20\%$ probability of $\beta^+$ decay and $40\%$ probability of electron capture. The mass of ${}^{64}_{29}$Cu is $63.92977$ amu while ${}^{64}_{30}$Zn is $63.92914$ amu and ${}^{64}_{28}$Ni is $63.92796$ amu.
\qn{3} What is the half life for electron capture(in hrs) ?
\qn{4} What is the Q (in MeV) value of $\beta^-$ decay?
\paraq{5}{6}
A point source $S$ emitting light of wavelength $\lambda=600$ nm is placed at a height equal to $600$ nm above a flat reflecting surface $AB$ (see figure). The intensity of the reflected light is $64\%$ of the incident intensity. Interference fringes are observed on a screen placed parallel to the reflecting surface at distance $2\lambda$ from it :-
\qfig{p1_phy_s3_para5_6_1}
\qn{5} The ratio of minimum to maximum intensity in the interference pattern observed on screen close to point '$P$' is :-
\qn{6} The radius (in nm) of the smallest bright ring formed on screen is :-
\sect{SECTION-II(ii)}
\qn{7} A stationary source emits sound of frequency $f_0=492$ Hz. The sound is reflected by a large car approaching the source with a speed of $2\ \mathrm{ms^{-1}}$. The reflected signal is received by the source and superposed with the original. What will be the beat frequency of the resulting signal in Hz ? (Given that the speed of sound in air is $330\ \mathrm{ms^{-1}}$ and the car reflects the sound at the frequency it has received).
\qn{8} In an $\alpha$-decay the kinetic energy of $\alpha$-particle is $48$ MeV and $Q$ value of the reaction is $50$ MeV. The mass number of the parent nucleus is $20n$ then find $n$ :- (Assume that daughter nucleus is in ground state)
\qn{9} An electron in an excited state of $\mathrm{Li^{2+}}$ ion has angular momentum $3h/2\pi$. The de Broglie wavelength of the electron in this state is $p\pi a_0$ (where $a_0$ is the Bohr radius). The value of $p$ is :-
\partsec{2}{CHEMISTRY}
\sect{SECTION-I(i)}
\qn{1} The emf of following three galvanic cells are found to be as $E_1$, $E_2$ and $E_3$ volts.\par
(i) $\mathrm{Zn\,|\,Zn^{2+}\,(0.01M)\,||\,Cu^{2+}\,(1M)\,|\,Cu}$\par
(ii) $\mathrm{Zn\,|\,Zn^{2+}\,(1M)\,||\,Cu^{2+}(0.01\ M)\,|\,Cu}$\par
(iii) $\mathrm{Zn\,|\,Zn^{2+}(1M)\,||\,Cu^{2+}\,(2M)\,|\,Cu}$\par
Choose the correct option :
\begin{qoptions}
\opt{A}{$E_1>E_2>E_3$}
\opt{B}{$E_3>E_2>E_1$}
\opt{C}{$E_1>E_3>E_2$}
\opt{D}{$E_2>E_3>E_1$}
\end{qoptions}
\qn{2} Which of the following is not reducing disaccharide ?
\begin{qoptions}
\opt{A}{Starch}
\opt{B}{Lactose}
\opt{C}{Maltose}
\opt{D}{Sucrose}
\end{qoptions}
\qn{3} Which of the following reaction can not give $\mathrm{PH_3}$ gas :-
\begin{qoptions}
\opt{A}{$\mathrm{PH_4I+KOH\to}$}
\opt{B}{$\mathrm{PCl_3+3H_2O\to}$}
\opt{C}{$\underset{\text{(white)}}{\mathrm{P_4}}+\mathrm{NaOH}\to$}
\opt{D}{$\mathrm{AlP+dil.\ HCl\to}$}
\end{qoptions}
\qn{4} A chemical which stimulates the secretion of pepsin is :
\begin{qoptions}
\opt{A}{Anti histamine}
\opt{B}{Cimetidine}
\opt{C}{Histamine}
\opt{D}{Zantac}
\end{qoptions}
\sect{SECTION-I(ii)}
\qn{5} Which of the following statements is/are correct ?
\begin{qoptions}
\opt{A}{Specific conductance is an intensive property while molar conductivity is an extensive property of solution.}
\opt{B}{The limiting equivalent conductivity of weak electrolyte can not be determined exactly by extrapolation of the plot of $\Lambda_{eq}$ against $\sqrt C$ [Given : $C$ = concentration of weak electrolyte]}
\opt{C}{The conductance of metals is due to the movements of free electrons.}
\opt{D}{Molar conductance of $0.01$ KCl aqueous solution is more than that of $0.1$ M KCl aqueous solution}
\end{qoptions}
\qn{6} The standard boiling point of benzene ($\mathrm{C_6H_6}$) is $80.1^\circ$C, then for given process $\mathrm{C_6H_6(\ell)\longrightarrow C_6H_6(g)}$. select the correct statement(s) for this process. Given phase diagram below is for benzene :\par
\qfig{p1_chem_s2_q06_1}
\begin{qoptions}
\opt{A}{$\Delta G$ is zero at standard boiling point.}
\opt{B}{$\Delta G^\circ$ is zero at standard boiling point.}
\opt{C}{$\Delta H^\circ$ is positive for given process.}
\opt{D}{$\Delta S^\circ$ is positive for given process}
\end{qoptions}
\qn{7} Given :\par
\qfig{p1_chem_s2_q07_1}
The compound(s), which on reaction with $\mathrm{HNO_3}$ will give the product having degree of rotation $[\alpha]_D=-52.7^\circ$ is (are)
\begin{qoptions}
\opt{A}{\qfigi{p1_chem_s2_q07_A}}
\opt{B}{\qfigi{p1_chem_s2_q07_B}}
\opt{C}{\qfigi{p1_chem_s2_q07_C}}
\opt{D}{\qfigi{p1_chem_s2_q07_D}}
\end{qoptions}
\qn{8} In which of the following reaction phenol obtained as product
\begin{qoptions}
\opt{A}{\qfigi{p1_chem_s2_q08_A}}
\opt{B}{\qfigi{p1_chem_s2_q08_B}}
\opt{C}{\qfigi{p1_chem_s2_q08_3}\ $\xrightarrow{\mathrm{KMnO_4}}$}
\opt{D}{\qfigi{p1_chem_s2_q08_D}\ $\xrightarrow{\text{Zn dust}}$}
\end{qoptions}
\qn{9} Which of the following statement(s) is/are correct?
\begin{qoptions}
\opt{A}{When NaOH is added to $\mathrm{NH_4Cl}$ then ammonia is evolved}
\opt{B}{$\mathrm{K_2CO_3}$ cannot be prepared by Solvay process}
\opt{C}{The alkali metals dissolve in liquid $\mathrm{NH_3}$ giving deep blue solution which are conducting in nature}
\opt{D}{$\mathrm{NaHCO_3}$ with $\mathrm{H_2SO_4}$ is used in fire extinguisher}
\end{qoptions}
\qn{10} Which of the following statements is/are CORRECT for $\mathrm{K_2MnO_4}$ ?
\begin{qoptions}
\opt{A}{It is prepared by $\mathrm{MnO_2}$ in acidic medium}
\opt{B}{It produce brown coloured compound on addition of water}
\opt{C}{It produce $\mathrm{KMnO_4}$ on addition of water}
\opt{D}{It produce purple colour compound with $\mathrm{O_3}$}
\end{qoptions}
\sect{SECTION-II(i)}
\paraq{1}{2}
Amount of heat evolved during complete combustion of liquid benzene can be calculated from the following data.\par
(i) $18$ gm of graphite on complete combustion evolves $591$ kJ heat in standard conditions\par
(ii) $\Delta H_f^\circ(\mathrm{H_2O,\ \ell})=-286$ kJ/mol\par
(iii) The heat of formation of liquid benzene is $50$ kJ /mole in standard conditions.
\qn{1} Heat of formation of $\mathrm{CO_2(g)}$ from above data (in kJ/mole) in standard conditions is -
\qn{2} Magnitude of heat evolved (in kJ) from combustion of $78$ gm benzene ($\ell$) in standard conditions is :
\paraq{3}{4}
Electrophilic aromatic substitution in phenols, the reaction that take place on the aromatic ring is electrophilic substitution reactions. The $-$OH group attached tot he benzene ring activates it towards electrophilic substitution. Also, it directs the incoming group to ortho and para positions in the ring as these positions become electron rich due to the mesomeric effect caused by $-$OH group.
\qn{3} \qfig{p1_chem_s3_q03_1}\par
Position of CHO group in A is $x$, no. of Br atoms in B is $y$ then $x+y=$ ?
\qn{4} An organic compound ($\mathrm{C_8H_{10}O_2}$) rotates plane-polarized light. It produces pink color with neutral $\mathrm{FeCl_3}$ solution. What is the total number of all the possible isomers for this compound?
\paraq{5}{6}
\qfig{p1_chem_s3_para5_6_1}
\qn{5} Find total number of atoms present in one plane of molecule $A$.
\qn{6} Calculate basicity of compound '$B$'.
\sect{SECTION-II(ii)}
\qn{7} The molar solubility of $\mathrm{Zn(OH)_2}$ in $C\times10^{-2}$ M NaOH solution is $2\times10^{-18}$ M. The value of $C$ is \underline{\hspace{1cm}}(Nearest integer)\par
(Given : The solubility product of $\mathrm{Zn(OH)_2}$ is $2\times10^{-20}$)
\qn{8} Find the number of compound(s) which are green colour .\par
$\mathrm{MnO_2}$, $\mathrm{K_2MnO_4}$, $\mathrm{Cr_2O_3}$, $\mathrm{Ni(OH)_2}$, $\mathrm{CuO}$, $\mathrm{CuCO_3{\cdot}Cu(OH)_2}$, $\mathrm{FeSO_4{\cdot}7H_2O}$
\qn{9} Find the number of ions which are soluble in excess of NaOH as well as excess of $\mathrm{NH_4OH}$\par
$\mathrm{Ag^+}$, $\mathrm{Pb^{+2}}$, $\mathrm{Al^{+3}}$, $\mathrm{Sn^{+2}}$, $\mathrm{Zn^{+2}}$, $\mathrm{Cd^{+2}}$.
\partsec{3}{MATHEMATICS}
\sect{SECTION-I(i)}
\qn{1} A bag contains $5$ white and $3$ black balls. $4$ balls are drawn one by one, their color noted and replaced in the bag. Probability that white and black balls are drawn alternately is
\begin{qoptions}
\opt{A}{$\dfrac{15^2}{2^{10}}$}
\opt{B}{$\dfrac{15^2}{2^{11}}$}
\opt{C}{$\dfrac{15^2}{2^{12}}$}
\opt{D}{$\dfrac{15^2}{2^{13}}$}
\end{qoptions}
\qn{2} If locus of point $P(z)$ which satisfies $\arg(3+z+\bar z+|z|i)=\dfrac\pi4$ is a part of conic $C$, then length of latus rectum of '$C$', is
\begin{qoptions}
\opt{A}{$1$}
\opt{B}{$\sqrt3$}
\opt{C}{$2\sqrt3$}
\opt{D}{$6$}
\end{qoptions}
\qn{3} The equation $x\ell nx+1=kx$ has -
\begin{qoptions}
\opt{A}{exactly one solution for $k=2$}
\opt{B}{exactly two solutions for $k=-2$}
\opt{C}{exactly one solution for $k=1$}
\opt{D}{exactly two solutions for $k=-1$}
\end{qoptions}
\qn{4} A vector in the plane of $\vec a=3\hat i+4\hat j$ and $\vec b=\hat i-\hat j$, whose projection on $\vec a$ is $7$ units and which is perpendicular to $\vec b$ is -
\begin{qoptions}
\opt{A}{$6\hat i+6\hat j$}
\opt{B}{$5\hat i+5\hat j$}
\opt{C}{$-8\hat i-8\hat j$}
\opt{D}{$-4\hat i-4\hat j$}
\end{qoptions}
\sect{SECTION-I(ii)}
\qn{5} If $\displaystyle\sum_{i=1}^{10}(x_i-5)=5$ and $\displaystyle\sum_{i=1}^{10}(x_i-5)^2=125$, variance of $x_1,x_2,\ldots\ldots x_{10}$ will be-
\begin{qoptions}
\opt{A}{$\dfrac94$}
\opt{B}{$\dfrac{25}{4}$}
\opt{C}{$\dfrac{49}{4}$}
\opt{D}{$\dfrac{81}{4}$}
\end{qoptions}
\qn{6} If $f(x)=\sqrt{x+\sqrt{x+\sqrt{x+\ldots\ldots}}}$, $x>0$, then
\begin{qoptions}
\opt{A}{$\dfrac{dy}{dx}=\dfrac{1}{2\sqrt{x+y}-1}$}
\opt{B}{$\dfrac{dy}{dx}=\dfrac{1}{2y-1}$}
\opt{C}{$\dfrac{dy}{dx}=\dfrac{1}{\sqrt{1+4x}}$}
\opt{D}{$\dfrac{dy}{dx}=\dfrac{1}{\sqrt{1+2x}}$}
\end{qoptions}
\qn{7} If $y(x)$ is the solution of $(2+y)\dfrac{dy}{dx}+(x-1)=0$ and $y(0)=0$ and the area bounded by the curve $y(x)$, $x$-axis and $y$-axis in first quadrant is $A$ then-
\begin{qoptions}
\opt{A}{$y(x)$ is circle of radius $5$.}
\opt{B}{$y(x)$ is a circle of radius $\sqrt5$}
\opt{C}{$\dfrac{5\pi}8-2<A<\dfrac{5\pi}{6}-2$}
\opt{D}{$\dfrac{5\pi}6-2<A<\dfrac{5\pi}{4}-2$}
\end{qoptions}
\qn{8} Given the equation of curve is $y=Ae^x\cos x+Be^x\sin x$, a differential equation of the given curve is formed, then which of the following is/are true-
\begin{qoptions}
\opt{A}{Order of differential equation is two}
\opt{B}{Degree of differential equation is one}
\opt{C}{Differential equation is linear}
\opt{D}{Differential equation is $\dfrac{d^2y}{dx^2}=2\left(y+\dfrac{dy}{dx}\right)$}
\end{qoptions}
\qn{9} Let $f$ be a twice differentiable function on $[0,2]$ such that $f(0)=0$, $f(1)=2$, $f(2)=4$, then
\begin{qoptions}
\opt{A}{$f'(\alpha)=2$ for some $a\in(0,1)$}
\opt{B}{$f''(x)=0\ \forall\ x\in(0,2)$}
\opt{C}{$f'(\beta)=2$ for some $\beta\in(1,2)$}
\opt{D}{$f''(\gamma)=0$ for some $\gamma\in(0,2)$}
\end{qoptions}
\qn{10} For $f(x)=e^{\int_0^x(t-3)(t-5)dt}$ which of following is/are correct ?
\begin{qoptions}
\opt{A}{Equation of normal at the point where curve cuts y-axis is $15y+x=0$}
\opt{B}{Equation of tangent at the point where curve cuts y-axis is $y=15x+1$}
\opt{C}{Number of integral values of $x$ at which function is strictly decreasing is one.}
\opt{D}{Number of integral values of $x$ at which function is strictly decreasing CANNOT be determined.}
\end{qoptions}
\sect{SECTION-II(i)}
\paraq{1}{2}
Let $A,B,C$ be $3$ exhaustive events of an experimental performance such that $P(A\cap B\cap C)=0$ and $P(A\cap B)=P(B\cap C)=P(C\cap A)=d\neq0$. If $P\left(A\bar B\cup\bar AB\right)=\dfrac{17}{24}$, $P\left(C\bar A\bar B\right)=\dfrac5{24}$ and $P\left(A\bar B\bar C\right)=\dfrac8{24}$. On the basis of above information, answer the following questions :
\qn{1} The probability of occurence of event $A$ given that $B$ has occured and $C$ has not occured , is-
\qn{2} The probability of occurence of event $A$ \& $B$ given that $C$ has not occured, is-
\paraq{3}{4}
Let $C_1$ and $C_2$ be the two curves on the complex plane defined as\par
$C_1:z+\bar z=2\,|z-1|$\par
$C_2:\arg(z+1+i)=\alpha$\par
where $\alpha$ belongs to the interval $(0,\pi)$ such that curves $C_1$ and $C_2$ touches each other at $P(z_0)$.
\qn{3} A particle starts from a point $P(z_0)$. It moves horizontally away from origin by $2$ units and then vertically away from origin by $3$ units to reach at a point $Q(z_1)$. If $z_1=x_1+iy_1$, then $(x_1+y_1)$ equals -
\qn{4} If $P(z_0)$ is rotated about origin through an angle '$2\alpha$' in clockwise direction, then the area bounded (in sq. unit) by $C_1$ and the line joining $P(z_0)$ and $Q(z_0')$ is (where $Q(z_0')$ is the new position $P(z_0)$ after rotation) sq. unit-
\paraq{5}{6}
Let $L$ be the line $\dfrac{x-6}{2}=\dfrac{y-12}{4}=\dfrac{z-8}{2}$. Let $A$ be the point of intersection of the plane $\pi_1:x+y+z=6$ \& $L$. Let $P$ be the point of intersection of the plane $\pi_2:x+y+z=10$ \& $L$. $AP$ is produced to $Q$ such that $AP=PQ$. Let $R$ be the image of $Q$ in $\pi_2$.
\qn{5} If the co-ordinates of $R$ is $(x_1,y_1,z_1)$ then the value of $(x_1+y_1+z_1)$ is-
\qn{6} Let $S$ be a point on the plane $\pi_2$ such that $PQSR$ is a parallelogram then its area is $\dfrac{\lambda\sqrt2}{3}$, then the value of '$\lambda$' is -
\sect{SECTION-II(ii)}
\qn{7} If $p,q,r$ \& $x,y,z$ are non-zero real numbers such that $px+qy+rz+\sqrt{(p^2+q^2+r^2)\cdot(x^2+y^2+z^2)}=0$. Then the value of $\dfrac{py}{qx}+\dfrac{qz}{ry}+\dfrac{rx}{pz}$ is equal to
\qn{8} Let $B$ be a skew symmetric matrix of order $3\times3$ with real entries. If $A=(I+B)(I-B)^{-1}$ where $I+B$, $I-B$ are non-singular matrices, then sum of values of $n$ satisfying the equation\par
{\small$\left|\sqrt[3]{n^2}\cdot A\right|-n\left|\sqrt[3]{5}\cdot\mathrm{adj}\,(\mathrm{adj}A)\right|+6|\mathrm{adj}A|^3=0$} is equal to (where $|A|$ denotes determinant value of $A$.)
\qn{9} If $\displaystyle\sum_{r=0}^{20}(2r+3^r)\,\left({}^{20}C_r\right)-\sum_{r=0}^{40}{}^{40}C_r=2^a.5$, then '$a$' is
\papersec{2}
\partsec{1}{PHYSICS}
\sect{SECTION-I(i)}
\qn{1} Moseley's law for $K_\alpha$ photon is given by $\sqrt\upsilon=a(Z-b)$ where $a$ is a universal constant and $b$ is a screening constant. Moseley's logic helps us assume $b=1$ for $K_\alpha$ photon. If there is a percentage error $=10^{-2}$ in the measurement of '$b$' due to actual orbital configuration of an atom, then the relative error in the measurement of $\upsilon$ if $Z=51$ is $p\times10^{-q}$. Then mark the correct options.
\begin{qoptions}
\opt{A}{Value of $p$ is $=4$}
\opt{B}{Value of $p$ is $=3$}
\opt{C}{Value of $q$ is $=7$}
\opt{D}{Value of $q$ is $=6$}
\end{qoptions}
\qn{2} Energy liberated in the de-excitation of hydrogen atom from $3^{\mathrm{rd}}$ level to $1^{\mathrm{st}}$ level falls on a photo-cathode. Later when the same photocathode is exposed to a spectrum of some unknown hydrogen like gas, excited to $2^{\mathrm{nd}}$ energy level, it is found that the de-Broglie wavelength of the fastest photoelectrons, now ejected has decreased by a factor of $3$. For this new gas, difference of energies of $2^{\mathrm{nd}}$ Lyman line and $1^{\mathrm{st}}$ Balmer line is found to be $3$ times the ionization potential of the hydrogen atom. Select the correct statement(s) :
\begin{qoptions}
\opt{A}{The gas is lithium.}
\opt{B}{The gas is helium.}
\opt{C}{The work function of photocathode is $8.5\ eV$.}
\opt{D}{The work function of photocathode is $5.5\ eV$.}
\end{qoptions}
\qn{3} A semicircular conducting ring of radius $R$ is placed in the $xy$ plane, as shown in the figure. A uniform magnetic field is set up along the $x$-axis. No emf, will be induced in the ring, if\par
\qfig{p2_phy_s1_q03_1}
\begin{qoptions}
\opt{A}{it moves along the $x$-axis.}
\opt{B}{it moves along the $y$-axis.}
\opt{C}{it moves along the $z$-axis.}
\opt{D}{it remains stationary.}
\end{qoptions}
\qn{4} A fish, $F$ in the pond, is at a depth of $0.8$ m from water surface and is moving vertically upwards with velocity $2\ \mathrm{ms^{-1}}$. At the same instant, a bird $B$ is at a height of $6$ m from water surface and is moving downwards with velocity $3\ \mathrm{ms^{-1}}$. At this instant both are on the same vertical lines as shown in the figure. Which of the following statement(s) is(are) correct?\par
\qfig{p2_phy_s1_q04_1}
\begin{qoptions}
\opt{A}{Height of $B$, observed by $F$ (from itself) is equal to $8.00$ m.}
\opt{B}{Depth of $F$, observed by $B$ (from itself) is equal to $6.60$ m.}
\opt{C}{Velocity of $B$, observed by $F$ (relative to itself) is equal to $5.00\ \mathrm{ms^{-1}}$.}
\opt{D}{Velocity of $F$, observed by $B$ (relative to itself) if equal to $4.50\ \mathrm{ms^{-1}}$.}
\end{qoptions}
\qn{5} A mass of $0.2$kg is attached to the lower end of a massless spring of force-constant $200$ N/m, the upper end of which is fixed to a rigid support. Which of the following statements is/are true?
\begin{qoptions}
\opt{A}{In equilibrium, the spring will be stretched by $1$cm.}
\opt{B}{If the mass is raised till the spring is unstretched state and then released, it will go down by $2$cm before moving upwards.}
\opt{C}{The frequency of oscillation will be nearly $5$ Hz.}
\opt{D}{If the system is taken to the moon, the frequency of oscillation will be the same as on the earth.}
\end{qoptions}
\qn{6} When the atomic number $A$ of the nucleus increases
\begin{qoptions}
\opt{A}{initially the neutron-proton ratio is constant $=1$}
\opt{B}{initially neutron-proton ratio increases and later decreases.}
\opt{C}{initially binding energy per nucleon increases and later decreases.}
\opt{D}{the binding energy per nucleon increases when the neutron-proton ratio increases.}
\end{qoptions}
\sect{SECTION-I(ii)}
\paraq{7}{8}
A cylindrical block of density $2$ gm/cc is suspended by a thin wire $3$m long. When the wire is vibrated, it produces fundamental frequency of $280$ Hz. The mass of wire is so small that tension in the wire can be assumed to be constant. Now the block is dipped in a liquid. It is found that the first overtone now is $280$ Hz.
\qfig{p2_phy_s2_para7_8_1}
\qn{7} What is the density of the liquid?
\begin{qoptions}
\opt{A}{$1$ gm/cc}
\opt{B}{$1.5$ gm/cc}
\opt{C}{$0.75$ gm/cc}
\opt{D}{$1.25$ gm/cc}
\end{qoptions}
\qn{8} If the block is dipped to $1/3$ of it's height in the liquid, what is the new fundamental frequency?
\begin{qoptions}
\opt{A}{$140$ Hz}
\opt{B}{$210$ Hz}
\opt{C}{$140\sqrt3$ Hz}
\opt{D}{$140\sqrt2$ Hz}
\end{qoptions}
\paraq{9}{10}
A monochromatic beam of light falls on Young's double slit experiment apparatus as shown in figure. A thin sheet of glass is inserted in front of lower slit $S_2$.
\qfig{p2_phy_s2_para9_10_1}
\qn{9} The central bright fringe can be obtained
\begin{qoptions}
\opt{A}{At $O$ only}
\opt{B}{At $O$ or below $O$ only}
\opt{C}{At $O$ or above $O$ only}
\opt{D}{Anywhere on the screen}
\end{qoptions}
\qn{10} If central bright fringe is obtained on screen at $O$ :-
\begin{qoptions}
\opt{A}{$(\mu-1)\,t=d\sin\theta$}
\opt{B}{$(\mu-1)\,t=d\cos\theta$}
\opt{C}{$(\mu-1)t+d\sin\theta=0$}
\opt{D}{$\dfrac{t}{\mu-1}=\dfrac{d}{\sin\theta}$}
\end{qoptions}
\sect{SECTION-II(i)}
\paraq{1}{2}
In the given circuit, reading of voltmeter $V_1$ is $100\sqrt3$ volt, value of $R$ is $10\Omega$ and value of $\dfrac1{\sqrt{LC}}$ is $\sqrt2\times100\pi$ rad/sec.
\qfig{p2_phy_s3_para1_2_1}
\qn{1} If reading of voltmeter $V_2$ is $100\sqrt n$ volts then what is value of $n$?
\qn{2} What is current in the circuit in amperes?
\paraq{3}{4}
At $t=0$, switch is closed. Capacitors are initially uncharged. Current through battery at $t=0^+$ is $(10\alpha)$ Ampere while that at $t\to\infty$ is $(10/\beta)$ Ampere.
\qfig{p2_phy_s3_para3_4_1}
\qn{3} The value of $\alpha$ is \underline{\hspace{1cm}}
\qn{4} Value of $\beta$ is \underline{\hspace{1cm}}
\paraq{5}{6}
Two charge particles each of mass '$m$', carrying charge $+q$ and connected with each other by a massless inextensible string of length $2L$ are describing circular path in the plane of paper, each with speed $v=\dfrac{qB_0L}{m}$ (where $B_0$ is constant) about their centre of mass in the region in which an uniform magnetic field $\vec B$ exists into the plane of paper as shown in figure. Neglect any effect of electrical \& gravitational forces.
\qfig{p2_phy_s3_para5_6_1}
\qn{5} The magnitude of the magnetic field such that no tension is developed in the string will be $\dfrac{B_0}{N}$, find $N$ :-
\qn{6} If the actual magnitude of magnetic field is half to that of calculated in part (i) then tension in the string is $\dfrac{q^2B_0^2L}{Nm}$, where $N$ is least integer. Find the value of $N$.
\sect{SECTION-II(ii)}
\qn{7} A parallel plate capacitor of plate separation $2$ mm is connected in an electric circuit having source voltage $400$ V. If the plate area is $60\ \mathrm{cm^2}$, then the value of displacement current for $10^{-6}$ s will be $x.yz\times10^{-2}$ A. Find $x+y$. ($\epsilon_0=8.85\times10^{-12}$SI units)
\qn{8} A particle of mass $2$ g and charge $1$mC is held at rest on a frictionless horizontal surface at a distance of $1$ m from a fixed charge of $1$ mC. If the particle is released it will be repelled. The speed of the particle (in m/s) when it is at distance of $10$ m from the fixed charge is:
\qn{9} When two deuterons (${}_1\mathrm{H}^2$) fuse to from a helium nucleus ${}_2\mathrm{He}^4$, $23.6$ MeV energy is released. Find the binding energy (in MeV) of helium if it is $1.1$ MeV for each nucleon of deuterium.
\partsec{2}{CHEMISTRY}
\sect{SECTION-I(i)}
\qn{1} Consider the reaction sequence\par
\qfig{p2_chem_s1_q01_1}
Select the correct statements :
\begin{qoptions}
\opt{A}{A and B can undergo oxidation by periodic acid.}
\opt{B}{C on esterification can produce cyclic ester.}
\opt{C}{Compound C can show hydrogen bonding.}
\opt{D}{Conversion of B to C involve migration of $\mathrm{CH_3}$ group.}
\end{qoptions}
\qn{2} Correct statement(s) about \qfigi{p2_chem_s1_q02_1} is/are
\begin{qoptions}
\opt{A}{$1$ mole $A$ librates $\dfrac32$ mole of $\mathrm{H_2}$ on treatment with Na.}
\opt{B}{$A$ gives positive test with $\mathrm{FeCl_3}$}
\opt{C}{$A$ gives positive test with $\mathrm{NaHCO_3}$}
\opt{D}{$A$ gives white ppt test with Tollen's reagent}
\end{qoptions}
\qn{3} Which statement(s) is/are true for the reaction.\par
$\mathrm{C+D\longrightarrow F+G}$\par
Experimental rate law for this reaction is found to be : Rate $=K[C]^2[D]^{1/2}$
\begin{qoptions}
\opt{A}{If the concentration of $D$ is increased to four times, then the reaction rate will double.}
\opt{B}{Reaction must be a complex reaction.}
\opt{C}{Multiple elementary steps are involved in the reaction mechanism.}
\opt{D}{If the concentrations of $C$ and $D$ both are doubled, then the reaction rate will double.}
\end{qoptions}
\qn{4} In which of the following pairs of solutions will osmotic presence of the solution be nearly equal? Assume complete ionization of given salts.
\begin{qoptions}
\opt{A}{$0.04$ M $\mathrm{K_4[Fe(CN)_6]}$ and $0.10$ M $\mathrm{FeSO_4}$}
\opt{B}{$0.10$ M $\mathrm{K_4[Fe(CN)_6]}$ and $0.05$ M $\mathrm{FeSO_4}$ $\mathrm{(NH_4)_2SO_4{\cdot}6H_2O}$}
\opt{C}{$0.20$ M NaCl and $0.10$ M $\mathrm{BaCl_2}$}
\opt{D}{$0.05$ M $\mathrm{FeSO_4}$ $\mathrm{(NH_4)_2\,SO_4{\cdot}6H_2O}$ and $0.05$ M $\mathrm{KCl{\cdot}MgCl_2{\cdot}6H_2O}$}
\end{qoptions}
\qn{5} Assuming that the complex Barium dichlorodioxalatocobaltate(III) is water soluble. Select the correct statement(s) for the above complex.
\begin{qoptions}
\opt{A}{It is diamagnetic.}
\opt{B}{When $1$ mole of complex solution is treated with excess of $\mathrm{Na_2SO_4}$, $2$ mole of $\mathrm{BaSO_4}$ is precipitated.}
\opt{C}{Complex ion is $\mathrm{d^2sp^3}$ hybrid.}
\opt{D}{Only one isomer of the complex is optically active.}
\end{qoptions}
\qn{6} When crystals of potassium chloride is heated with $\mathrm{MnO_2}$ and conc. $\mathrm{H_2SO_4}$ then
\begin{qoptions}
\opt{A}{$\mathrm{Cl_2}$ is produced}
\opt{B}{Its redox change}
\opt{C}{Resulting mixture after dilution by water gives buff colour ppt. with $\mathrm{H_2S}$}
\opt{D}{Resulting mixture after dilution by water gives ppt. with excess KOH.}
\end{qoptions}
\sect{SECTION-I(ii)}
\paraq{7}{8}
The $-\mathrm{NO_2}$ group in an aromatic ring deactivates the ortho and para positions for an electrophilic attack. When $-\mathrm{NO_2}$ group is present at ortho or para positions of a leaving group (Nucleofuge) it activates the ring for nucleophilic attack. The reduction of $-\mathrm{NO_2}$ group by metal in acid causes its reduction to $-\mathrm{NH_2}$ group and then the ring becomes strongly activated for a electrophilic attack. The strong activating group ($-\mathrm{NH_2}$) is converted into moderately activating group ($-\mathrm{NH{-}C(=O){-}CH_3}$, i.e.\ $-$NHAc) by it's acylation with $\mathrm{CH_3{-}CO{-}Cl}$. Deacylation is carried out by hydrolysis with $\mathrm{H_3O^+}$ or $\mathrm{OH^-}$. The ring alkylation by using $\mathrm{RX/AlX_3}$ is not possible in presence of $-\mathrm{NO_2}$ or $-\mathrm{NH_2}$ group but is possible in presence of $-$NHAc group.\par
\qfig{p2_chem_s2_para7_8_4}
\qn{7} The product ($\mathrm{P_1}$) is :
\begin{qoptions}
\opt{A}{\qfigi{p2_chem_s2_q07_A}}
\opt{B}{\qfigi{p2_chem_s2_q07_B}}
\opt{C}{\qfigi{p2_chem_s2_q07_C}}
\opt{D}{\qfigi{p2_chem_s2_q07_D}}
\end{qoptions}
\qn{8} The product ($\mathrm{P_2}$) is :
\begin{qoptions}
\opt{A}{\qfigi{p2_chem_s2_q08_A}}
\opt{B}{\qfigi{p2_chem_s2_q08_B}}
\opt{C}{\qfigi{p2_chem_s2_q08_C}}
\opt{D}{\qfigi{p2_chem_s2_q08_D}}
\end{qoptions}
\paraq{9}{10}
Upon heating $\mathrm{KClO_3}$ in the presence of catalytic amount of $\mathrm{MnO_2}$, a gas $W$ is formed. Excess amount of $W$ reacts with white phosphorus to give $X$. The reaction of $X$ with pure $\mathrm{HNO_3}$ gives $Y$ and $Z$.
\qn{9} $Y$ and $Z$ are , respectively
\begin{qoptions}
\opt{A}{$\mathrm{N_2O_4}$ and $\mathrm{H_3PO_3}$}
\opt{B}{$\mathrm{N_2O_4}$ and $\mathrm{HPO_3}$}
\opt{C}{$\mathrm{N_2O_5}$ and $\mathrm{HPO_3}$}
\opt{D}{$\mathrm{N_2O_3}$ and $\mathrm{H_3PO_3}$}
\end{qoptions}
\qn{10} $W$ and $X$ are, respectively --
\begin{qoptions}
\opt{A}{$\mathrm{O_3}$ and $\mathrm{P_4O_6}$}
\opt{B}{$\mathrm{O_2}$ and $\mathrm{P_4O_{10}}$}
\opt{C}{$\mathrm{O_3}$ and $\mathrm{P_4O_{10}}$}
\opt{D}{$\mathrm{O_2}$ and $\mathrm{P_4O_6}$}
\end{qoptions}
\sect{SECTION-II(i)}
\paraq{1}{2}
Consider the following reactions\par
$\mathrm{D-(+)-glucose}\xrightarrow{\mathrm{HNO_3}}[\mathrm P]$\par
$\mathrm{D-(+)-glucose}\xrightarrow{\mathrm{Br_2/H_2O}}[\mathrm Q]$\par
The number of stereo isomers possible for $P$ and $Q$ (including $P$ \& $Q$), respectively, are $x$ and $y$.
\qn{1} The value of $x$ is \underline{\hspace{0.8cm}}.
\qn{2} The value of $y$ is \underline{\hspace{0.8cm}}.
\paraq{3}{4}
The external pressure $2$ atm is applied on frictionless movable piston, fitted in a vessel containing $100$g of $X(l)$ at $450$ K. Now heat is supplied keeping pressure constant till $40$g of $X$ is evaporated to form $X(g)$ at $500$ K(Boiling point). Assume vapour of $X(l)$ behaves like an ideal gas.
\qfig{p2_chem_s3_para3_4_1}
\textbf{Given :} Molar heat capacity of $X(l)=60$ J/mol. K; $\Delta H_{\text{vaporisation}}$ of $X(\ell)=30$ kJ/mol; $R=8.3$ J/mol-K ; Molar mass of $X=20$ g/mol.
\qn{3} Calculate change in internal ($\Delta U$) energy in kJ for overall process.
\qn{4} Calculate the magnitude of work done (in kJ) during vapourization of $40$ g $X(\ell)$ given above.
\paraq{5}{6}
The distribution of electrons among various molecular orbitals is called the electronic configuration of the molecule which provides us the following very important informations about the molecule.\par
\textbf{(A) Stability of molecule} : The molecule is stable if number of bonding molecular orbital electrons ($N_b$) is greater than the number of antibonding molecular orbital electrons ($N_a$).\par
\textbf{(B) Bond order} : Bond order $=(N_b-N_a)/2$\par
A positive bond order means a stable molecule while a negative or zero bond order means an unstable molecule.\par
\textbf{(C) Nature of bond} : Bond order $1$, $2$ or $3$ corresponds to single, double or triple bonds respectively.\par
\textbf{(D) Bond length :} Bond length decreases as bond order increases.\par
\textbf{(E) Magnetic nature} : Molecular orbitals in a molecule are doubly occupied, the substance is diamagnetic and it one or more molecular orbitals are singly occupied it is paramagnetic.
\qn{5} Find the no. of species having fractional bond order?\par
$\mathrm{N_2^+}$, $\mathrm{N_2^-}$, $\mathrm{O_2}$, $\mathrm{O_2^+}$, $\mathrm{F_2}$, $\mathrm{B_2^+}$, $\mathrm{C_2^+}$, $\mathrm{CN^-}$, $\mathrm{NO^+}$
\qn{6} Assuming there is no $2s{-}2p$ mixing, the number of paramagnetic species amongst the following :\par
$\mathrm{H_2}$, $\mathrm{Li_2}$, $\mathrm{Li_2^{\oplus}}$, $\mathrm{B_2}$, $\mathrm{B_2^+}$, $\mathrm{C_2}$, $\mathrm{N_2}$, $\mathrm{N_2^{2\oplus}}$, $\mathrm{O_2}$, $\mathrm{F_2^{\oplus}}$
\sect{SECTION-II(ii)}
\qn{7} How many of the following statements are True ?\par
(a) Dacron is condensed copolymer\par
(b) Butan-1-ol is more soluble in water than propan-1-ol\par
(c) Lactose is reducing sugar but it gives negative Barfoed test\par
(d) $\mathrm{CHI_3}$ has antiseptic properties while $\mathrm{CHCl_3}$ has aneasthetic properties\par
(e) p-nitrophenol is more acidic than o-nitrophenol\par
(f) Carboxylic acid having $\alpha$-H shows HVZ reaction\par
(g) $\mathrm{CH_3{-}Br}$ has higher rate of $\mathrm{S_N2}$ than \qfigi{p2_chem_s4_q07_1}\par
(h) $1,1$-dichlorobut-$2$-ene can't show geometrical isomerism\par
(i) Secondary alkyl amine is less basic than tertiary alkyl amine in aqueous medium
\qn{8} A crystal of lead(II) sulphide has NaCl-type structure. In this crystal the shortest distance between $\mathrm{Pb^{2\oplus}}$ ion and $\mathrm{S^{2\ominus}}$ ion is $300$ pm. What is the nearest distance (in pm) between nearest tetrahedral and octahedral voids in PbS lattice.
\qn{9} A dilute solution of sulphuric acid is electrolysed using a current of $0.10$ A for $2$ hours to produce hydrogen and oxygen gas. The total volume of gases produced at STP is \underline{\hspace{1cm}} $\mathrm{cm^3}$. (Nearest integer)\par
[Given : Faraday constant $F=96500\ \mathrm{C\,mol^{-1}}$ at STP, molar volume of an ideal gas is $22.7\ \mathrm{L\,mol^{-1}}$]
\partsec{3}{MATHEMATICS}
\sect{SECTION-I(i)}
\qn{1} $A$ and $B$ are two independent events of an experiment, such that $P(A)=\dfrac13$, $P(B)=\dfrac14$, then
\begin{qoptions}
\opt{A}{$P(A\cup B)=\dfrac12$}
\opt{B}{$P(A\cup\bar B)=\dfrac56$}
\opt{C}{$P(B\cup\bar A)=\dfrac34$}
\opt{D}{$P(A\cap B)=\dfrac1{12}$}
\end{qoptions}
\qn{2} Let $A=\begin{bmatrix}1&2&3\\10&20&30\\11&22&k\end{bmatrix}$ and $v=\begin{bmatrix}x\\y\\z\end{bmatrix}$, where $k$ is a constant and $x,y,z$ are variables. Then which of the following statement is/are True ?
\begin{qoptions}
\opt{A}{Regardless of the value of $k$, the matrix $A$ is not invertible, i.e., there is no $3\times3$ matrix $B$ such that $BA=I$ (where $I$ is $3\times3$ identity matrix)}
\opt{B}{There is a unique $k$ such that determinant of $A$ is $0$.}
\opt{C}{The set of solutions $(x,y,z)$ of the matrix equation $Av=\begin{bmatrix}0\\0\\0\end{bmatrix}$ is either a line or a plane containing the origin.}
\opt{D}{If the equation $Av=\begin{bmatrix}p\\q\\r\end{bmatrix}$ has a solution, then it must be true that $q=10p$.}
\end{qoptions}
\qn{3} Two circles of radii $a$ and $1$ touch each other externally and are inscribed in the region bounded by $y=-\sqrt{4-x^2}$ and $x$-axis, then
\begin{qoptions}
\opt{A}{$100(a+1)=125$}
\opt{B}{Number of such pairs of circles is $2$}
\opt{C}{$a=1/2$}
\opt{D}{Points of contact of circle with radius $a$ and $y=-\sqrt{4-x^2}$ are $\left(\dfrac{\pm4\sqrt2}{3},-\dfrac23\right)$}
\end{qoptions}
\qn{4} If $\vec p,\vec q,\vec r$ are three non zero coplanar vectors then $(\vec p+\vec q+\vec r)\cdot((\vec p+\vec q)\times(\vec p+\vec r))$ equals to
\begin{qoptions}
\opt{A}{$0$}
\opt{B}{$[\vec p\ \vec q\ \vec r]$}
\opt{C}{$2\,[\vec p\ \vec q\ \vec r]$}
\opt{D}{$-[\vec p\ \vec q\ \vec r]$}
\end{qoptions}
\qn{5} Two perpendicular unit vectors $\vec a$ and $\vec b$ are such that $[\vec r\ \vec a\ \vec b]=\dfrac54$, $\vec r\cdot(3\vec a+2\vec b)=0$ and $\displaystyle\int_{-2\vec r\cdot\vec a}^{\frac{-4}{3}\vec r\cdot\vec b}\dfrac{x+1}{x^2+1}dx=\dfrac\pi2$. Which of the following is(are) correct ?
\begin{qoptions}
\opt{A}{$|\vec r|^2=\dfrac{19}8$}
\opt{B}{$|\vec r|=\dfrac74$}
\opt{C}{$\vec r=\dfrac{\vec a}2-\dfrac{3\vec b}4+\dfrac54(\vec a\times\vec b)$}
\opt{D}{$\vec r=\dfrac{-\vec a}2+\dfrac{3\vec b}2+\dfrac54(\vec a\times\vec b)$}
\end{qoptions}
\qn{6} Let $\Delta PQR$ be a triangle. Let $\vec a=\overrightarrow{QR}$, $\vec b=\overrightarrow{RP}$ and $\vec c=\overrightarrow{RQ}$. If $|\vec a|=12$, $|\vec b|=4\sqrt3$ and $\vec b\cdot\vec c=24$, then which of the following is (are) true?
\begin{qoptions}
\opt{A}{$\dfrac{|\vec c|^2}{2}-|\vec a|=12$}
\opt{B}{$\dfrac{|\vec c|^2}{2}+|\vec a|=30$}
\opt{C}{$|\vec a\times\vec b+\vec c\times\vec a|=48\sqrt3$}
\opt{D}{$\vec a\cdot\vec b=-72$}
\end{qoptions}
\sect{SECTION-I(ii)}
\paraq{7}{8}
Let $f(x)$ is a continous and twice differentiable function for all $x\in R$, and\par
$f'(c_1)=0=f'(c_2)$, $f(c_1)=5$, $f(c_2)=0$\par
$f''(c_1)\,f''(c_2)<0$ ($c_1<c_2$)
\qn{7} If $f''(c_1)-f''(c_2)<0$, then minimum number of root of $f'(x)=0$ in $[c_1-1,\,c_2+1]$ is :
\begin{qoptions}
\opt{A}{$1$}
\opt{B}{$2$}
\opt{C}{$3$}
\opt{D}{None of these}
\end{qoptions}
\qn{8} If $f''(c_1)-f''(c_2)>0$ , then minimum number of root of $f'(x)=0$ in $[c_1-1,\,c_2+1]$ is :
\begin{qoptions}
\opt{A}{$2$}
\opt{B}{$3$}
\opt{C}{$4$}
\opt{D}{None of these}
\end{qoptions}
\paraq{9}{10}
Let $f(x)=\begin{cases}[x],&-3\leqslant x<-1\\x^2-4,&-1\leqslant x\leqslant3\end{cases}$ and $g(x)=f(|x|)+|f(x)|$ where $[\,.\,]$ denotes G.I.F.
\qn{9} The number of points where $|f(x)|$ is non differentiable are :
\begin{qoptions}
\opt{A}{$2$}
\opt{B}{$3$}
\opt{C}{$4$}
\opt{D}{$5$}
\end{qoptions}
\qn{10} The number of points where $g(x)$ is non differentiable are :
\begin{qoptions}
\opt{A}{$2$}
\opt{B}{$3$}
\opt{C}{$4$}
\opt{D}{$5$}
\end{qoptions}
\sect{SECTION-II(i)}
\paraq{1}{2}
Let $y=f(x)$ be a curve $C_1$ passing through $(2,2)$ and $\left(8,\dfrac12\right)$ and satisfying the differential equation $y\left(\dfrac{d^2y}{dx^2}\right)=2\left(\dfrac{dy}{dx}\right)^2$.
\qn{1} If shortest distance between $C_1$ and $C_2:x^2+y^2=4$ is $(\sqrt a-b)$ where $a,b\in N$, then $(a+b)$ is equal to
\qn{2} Let $S$ be the area bounded by curve $C_1$, $x$-axis, $x=1$ and $x=4$, then the value of $e^S$ is
\paraq{3}{4}
Consider $f(x)=\sin^{-1}x+\tan^{-1}x+\sqrt{x+1}$ \& let $g(x)$ be its inverse.
\qn{3} Number of integral values of $k$ such that $f(x)=k$ has atleast one real solution is equal to -
\qn{4} Number of distinct solutions of the equation $f(|x|)=1$ is equal to -
\paraq{5}{6}
Consider two planes $P_1:2x-3y+z=-5$ and $P_2:x-4y-14z=5$. If a line $L$ perpendicular to $P_1$ whose one point $A(-7,y_1,z_1)$ lies on plane $P_1$ and point $B(x_2,y_2,z_2)$ and $C(x_3,y_3,z_3)$ lies on plane $P_2$ as well as on line $L$.
\qn{5} Let equation of line $L$ be $\dfrac{x+7}{a}=\dfrac{y-y_1}{-b}=\dfrac{z-z_1}{c}$ where $a,b,c\in N$ then least value of $a+b+c$ is
\qn{6} Value of $\dfrac{z_3-z_2}{y_3-y_2}+\dfrac{x_3-x_2}{z_3-z_2}$ is
\sect{SECTION-II(ii)}
\qn{7} If $f(x)=(x^3-1)|x-1|+\sin(\pi x)|x-2|+\left|\,||x|-3|-5\right|$ and $\lambda$ be the number of points in which $f(x)$ in non-differentiable then last two digit of the number $\lambda^{32}$ is :
\qn{8} If $f(x)=x^2\left(\displaystyle\int_1^{x^2}f(t)dt\right)$, then $f(1)+f'(1)$ is
\qn{9} Let $A$ be a non-singular matrix of order $3$ satisfying the equation\par
$(A^3-4A^2)(A^2+4A+16\,I)=0$.\par
If $N=|\mathrm{adj}.(\mathrm{adj}.2A)|$ then $\log_2N$ is :
\end{multicols}
\end{document}
